\documentclass[11pt]{article}

%-------------------------
% Packages
%-------------------------
\usepackage[a4paper,margin=1in]{geometry}
\usepackage{xcolor}
\usepackage{tikz}
\usepackage[most]{tcolorbox}
\usepackage{mathtools}
\usepackage{amsthm, amsmath, amssymb}
\usepackage{enumitem}
\usepackage{hyperref}
\usepackage[nameinlink,noabbrev]{cleveref}
\usepackage{titling} 
\usepackage{fontspec}
\usepackage{unicode-math}


%-------------------------
% Fonts
%-------------------------
\setmainfont{JetBrains Mono}
\setsansfont{JetBrains Mono}
\setmonofont{JetBrains Mono}


\setmathfont[
    Scale=1.25
]{latinmodern-math.otf}

\setmathfont[
  range=\mathup/{latin, Latin, 0-9},
  range=\mathit/{latin, Latin},
  range=\mathbfit/{latin, Latin},
  Scale=1.5
]{JetBrains Mono Bold} 

\setmathfont[range=\mathbb, Scale=MatchUppercase]{latinmodern-math.otf}

%-------------------------
% Paragraph spacing
%-------------------------
\setlength{\parindent}{0pt}
\setlength{\parskip}{0.7em}


%-------------------------
% Dark mode colors
%-------------------------
\definecolor{bgcolor}{HTML}{333333}
\definecolor{textcolor}{HTML}{FAFAFA}
\definecolor{defcolor}{HTML}{E86873}
\definecolor{thmcolor}{HTML}{0A9396}
\definecolor{lemcolor}{HTML}{94D2BD}
\definecolor{corcolor}{HTML}{9B4AF7}
\definecolor{probcolor}{HTML}{EE9B00}
\definecolor{excolor}{HTML}{21E933}
\definecolor{propcolor}{HTML}{EE9B20}

%-------------------------
% Background and text color
%-------------------------
\pagecolor{bgcolor}
\color{textcolor}

%-------------------------
% Paragraph spacing
%-------------------------
\setlength{\parindent}{0pt}
\setlength{\parskip}{0.7em}

%-------------------------
% Base tcolorbox settings (Removed 'attach title to upper')
%-------------------------
\tcbset{
  enhanced,
  colback=bgcolor,
  colframe=thmcolor,
  coltext=white,
  coltitle=white,
  fonttitle=\bfseries,
  boxrule=0.7pt,
  left=1em,
  right=1em,
  top=0.7em,
  bottom=0.7em,
  before skip=10pt,
  after skip=10pt,
  title after break,
  breakable
}

%-------------------------
% Standard Theorem Environments (using amsthm)
% Syntax: \begin{thm}[Optional Title]\label{thm:label} ...
%-------------------------
\newtheorem{thm}{Theorem}[section]
\newtheorem*{thm*}{Theorem}
\newtheorem{defn}{Definition}[section]
\newtheorem*{defn*}{Definition}
\newtheorem{lem}{Lemma}[section]
\newtheorem*{lem*}{Lemma}
\newtheorem{cor}{Corollary}[section]
\newtheorem*{cor*}{Corollary}
\newtheorem{prob}{Problem}[section]
\newtheorem*{prob*}{Problem}
\newtheorem{ex}{Example}[section]
\newtheorem*{ex*}{Example}
\newtheorem{prop}{Proposition}[section]
\newtheorem*{prop*}{Proposition}

%-------------------------
% Tcolorbox environments (applies style to amsthm environments)
% Removed 'attach title to upper' from all individual settings.
%-------------------------
\tcolorboxenvironment{thm}{
  colframe=thmcolor,
  colback=thmcolor!15!bgcolor,
  fonttitle=\bfseries,
  coltitle=white,
  breakable,
  title after break,
}

\tcolorboxenvironment{thm*}{
  colframe=thmcolor,
  colback=thmcolor!15!bgcolor,
  fonttitle=\bfseries,
  coltitle=white,
  breakable,
  title after break,
}

\tcolorboxenvironment{defn}{
  colframe=defcolor,
  colback=defcolor!15!bgcolor,
  fonttitle=\bfseries,
  coltitle=white,
  breakable,
  title after break,
}

\tcolorboxenvironment{defn*}{
  colframe=defcolor,
  colback=defcolor!15!bgcolor,
  fonttitle=\bfseries,
  coltitle=white,
  breakable,
  title after break,
}

\tcolorboxenvironment{lem}{
  colframe=lemcolor,
  colback=lemcolor!15!bgcolor,
  fonttitle=\bfseries,
  coltitle=white,
  breakable,
  title after break,
}

\tcolorboxenvironment{lem*}{
  colframe=lemcolor,
  colback=lemcolor!15!bgcolor,
  fonttitle=\bfseries,
  coltitle=white,
  breakable,
  title after break,
}

\tcolorboxenvironment{cor}{
  colframe=corcolor,
  colback=corcolor!15!bgcolor,
  fonttitle=\bfseries,
  coltitle=white,
  breakable,
  title after break,
}

\tcolorboxenvironment{cor*}{
  colframe=corcolor,
  colback=corcolor!15!bgcolor,
  fonttitle=\bfseries,
  coltitle=white,
  breakable,
  title after break,
}

\tcolorboxenvironment{prob}{
  colframe=probcolor,
  colback=probcolor!15!bgcolor,
  fonttitle=\bfseries,
  coltitle=white,
  breakable,
  title after break,
}

\tcolorboxenvironment{prob*}{
  colframe=probcolor,
  colback=probcolor!15!bgcolor,
  fonttitle=\bfseries,
  coltitle=white,
  breakable,
  title after break,
}

\tcolorboxenvironment{ex}{
  colframe=excolor,
  colback=excolor!15!bgcolor,
  fonttitle=\bfseries,
  coltitle=white,
  breakable,
  title after break,
}

\tcolorboxenvironment{ex*}{
  colframe=excolor,
  colback=excolor!15!bgcolor,
  fonttitle=\bfseries,
  coltitle=white,
  breakable,
  title after break,
}

\tcolorboxenvironment{prop}{
    colframe=propcolor, 
    colback=propcolor!15!bgcolor,
    fonttitle=\bfseries,
    coltitle=white,
    breakable,
    title after break
}

\tcolorboxenvironment{prop*}{
    colframe=propcolor, 
    colback=propcolor!15!bgcolor,
    fonttitle=\bfseries,
    coltitle=white,
    breakable,
    title after break
}

%-------------------------
% Cleveref settings (Using short environment names for correct display)
%-------------------------
\crefname{thm}{theorem}{theorems}
\Crefname{thm}{Theorem}{Theorems}

\crefname{defn}{definition}{definitions}
\Crefname{defn}{Definition}{Definitions}

\crefname{lem}{lemma}{lemmas}
\Crefname{lem}{Lemma}{Lemmas}

\crefname{cor}{corollary}{corollaries}
\Crefname{cor}{Corollary}{Corollaries}

\crefname{prob}{problem}{problems}
\Crefname{prob}{Problem}{Problems}

\crefname{ex}{example}{examples}
\Crefname{ex}{Example}{Examples}

\crefname{prop}{proposition}{propositions}
\Crefname{prop}{Proposition}{Propositions}

\crefname{thm*}{theorem}{theorems}
\Crefname{thm*}{Theorem}{Theorems}

\crefname{defn*}{definition}{definitions}
\Crefname{defn*}{Definition}{Definitions}

\crefname{lem*}{lemma}{lemmas}
\Crefname{lem*}{Lemma}{Lemmas}

\crefname{cor*}{corollary}{corollaries}
\Crefname{cor*}{Corollary}{Corollaries}

\crefname{prob*}{problem}{problems}
\Crefname{prob*}{Problem}{Problems}

\crefname{ex*}{example}{examples}
\Crefname{ex*}{Example}{Examples}

\crefname{prop*}{proposition}{propositions}
\Crefname{prop*}{Proposition}{Propositions}


\usepackage{tikz}
\usepackage{tikz-cd}

\DeclareMathOperator{\id}{\operatorname{id}}

\title{\huge{Assignment 4}}
\author{\LARGE{Thobias Høivik | Remote group}}
\date{for October 6, 2026}

\begin{document}
\maketitle

\begin{prob*}[Munkres 23.10]
    Let $\{X_\alpha\}_{\alpha \in J}$ 
    be an indexed family of 
    connected spaces; 
    let $X$ be the product space 
    \[
        X = \prod_{\alpha \in J} X_\alpha. 
    \]

    Let $a = (a_\alpha)$ be a fixed point 
    of $X$. 

    \begin{enumerate}[label=(\alph*)]
        \item Given any finite subset 
            $K$ of $J$, let $X_K$ 
            denote the subspace of $X$ 
            consisting of all 
            points $x = (x_\alpha)$ 
            such that $x_\alpha = a_\alpha$ 
            for $\alpha \not\in K$. 
            Show that $X_K$ is connected. 
        \item Show that the union $Y$
            of the spaces $X_K$ is connected. 
        \item Show that $X$ equals the 
            closure of $Y$. Conclude 
            that $X$ is connected. 
    \end{enumerate}
\end{prob*}
\begin{proof}[Proof of (a)]
    $K\subseteq J$ is a finite subset, say 
    \[
        K = \{\alpha_1,\ldots,\alpha_n\}. 
    \]
    Then $x = (x_\alpha) \in X_K$ coincides with 
    $a$ on all, but finitely many $x_\alpha$, namely 
    the indices 
    \[
        \{x_{\alpha_i}\}_{i\in[n]}. 
    \]
    More importantly, $x,y \in X_K$ coincide on all 
    but these finitely many indices. 

    The most straightforward approach for 
    connectedness is to show that 
    \[
        X_K \cong \prod_{i\in [n]} X_{\alpha_i}. 
    \]
    Connectedness is preserved under homeomorphism.

    Let $\varphi : X_K \to \prod_{i\in[n]} X_\alpha$ by 
    be the map that sends 
    \[
        (x_\alpha) \to (x_{\alpha_1}, x_{\alpha_2}, \ldots, 
        x_{\alpha_n}).  
    \]
    
    Since, for all $x,y \in X_K$, they coincide 
    on all but these finitely many indices. 

    Thus there is a 
    natural inverse $\varphi^{-1} : \prod_{i\in[n]} \to X_K$ 
    which sends $(x_{\alpha_{1}}, \ldots, x_{\alpha_n})$ 
    to the $x = (x_\alpha) \in X_K$ for 
    which the finitely many nonfixed indices are 
    $x_{\alpha_1}, x_{\alpha_2}, \ldots, x_{\alpha_n}$. 

    For rigor, define the injection 
    $\iota : [n] \to J$ 
    which sends $i$ to the index in $J$ corresponding 
    to $\alpha_i$. 

    Then $\varphi^{-1}$ is the function mapping 
    \[
        (x_{\alpha_1}, \ldots, x_{\alpha_n}) 
        \mapsto y = (y_\alpha)_{\alpha \in J} 
    \]
    where $y_{\iota_i} = x_{\alpha_i}$ and 
    $y_\alpha = a_\alpha$ otherwise. 

    This map is well-defined since 
    $y = \varphi^{-1}(x_{\alpha_1}, \ldots, x_{\alpha_n})$ 
    coincides with $a$ on all, but the finitely 
    many indices. 

    It is not terribly hard to see 
    that $\varphi^{-1}\circ \varphi \equiv \id_{K_X}$ 
    and that $\varphi \circ \varphi^{-1} \equiv \id_{\prod X_\alpha}$. 

    It remains to prove that these maps are continuous.

    Let $U \subseteq \prod_{i\in[n]} X_{\alpha_i}$ be 
    an open set. 

    Then, 
    \[
        U = \bigcup_{\beta} 
        \prod_{i=1}^n U_i, \ U_i \ \text{open in} \ X_{\alpha_i}, 
    \]
    which is annoying. 

    Luckily, it suffices to show the 
    preimage of 
    basis elements are open in the domain. 

    The basiselements of the topology on $\prod_{i\in [n]} 
    X_{\alpha_i}$ are 
    \[
        \prod_{i=1}^n U_i, \ U_i \ \text{open in}\ X_{\alpha_i},  
    \]
    and the ones in $X_K$ are 
    \[
        \prod_{\alpha} U_\alpha, 
    \]
    where $U_\alpha$ are open in $X_K$ and 
    only finitely many $U_\alpha \neq X_K$. 

    For a basis element $\prod U_i$ of $\prod_i X_{\alpha_i}$, 
    we have 
    \[
        \varphi^{-1}\left(\prod_{i=1}^n U_i\right) 
        = \prod_{\alpha} U_\alpha
    \]
    where only those $U_\beta$ for which 
\end{proof} 



\end{document}
